\titledquestion{Multiple Choices}

Each question has \textbf{one or more} correct answer(s). Select all the correct answer(s). Each part is worth 2 points: 2 for exactly the correct set, 1 for a non-empty proper subset of it, and 0 for an empty selection or any incorrect selection.

Write your answers in the following table.

\begin{table}[htbp]
	\centering
	\begin{tabular}{|p{1.6cm}|p{1.6cm}|p{1.6cm}|p{1.6cm}|p{1.6cm}|p{1.6cm}|}
		\hline
		(a) & (b) & (c) & (d) & (e) & (f) \\
		\hline
		C & AD & ABD & D & AB & AC \\
		\hline
	\end{tabular}
\end{table}

\begin{parts}
	\part[2] Assume we insert one element before the element at index \(i(0\leqslant i\leqslant n-1)\) in an array of length \(n\) and get a new array of length \(n+1\). Assume the array has a spare slot at the end, keeps its starting address, and preserves element order. Do not count writing the new element. What is the minimum number of moves that the elements in the array need in total?

	\begin{choices}
		\choice \(0\)
		\choice \(i\)
		\choice \(n-i\)
		\choice \(n-i-1\)
	\end{choices}

	\part[2] Which of the following statements about arrays and linked-lists are true?

	\begin{choices}
		\choice Gaining access to the \(k\)-th element in an array takes constant time.
		\choice Gaining access to the \(k\)-th element in a linked-list takes constant time.
		\choice Erasing the \(k\)-th element in an array takes constant time.
		\choice With access to the \(k\)-th element, inserting an element after the \(k\)-th element in a linked-list takes constant time.
	\end{choices}

    \part[2] Linear search checks an array of \(n\) elements from left to right and stops when it finds the target, or after checking every element. Count element comparisons. Which statements are true?
    \begin{choices}
        \choice Its best-case running time is \(\Theta(1)\).
        \choice Its worst-case running time is \(\Theta(n)\).
        \choice An \(O(n)\) upper bound implies that every search takes \(\Theta(n)\) time.
        \choice If the target occurs exactly once and is equally likely to be at any index, the expected number of comparisons is \((n+1)/2\).
    \end{choices}

    \part[2] A queue uses a circular array indexed from \(0\) to \(n-1\), with \(n\ge2\). A separate size counter distinguishes full from empty, so all \(n\) slots can be used. If the queue is full and its first element is at index \(m\), where is its last element? Interpret each option modulo \(n\).

	\begin{choices}
		\choice \(0\)
		\choice \(m\)
		\choice \(n-1\)
		\choice \(m-1\)
	\end{choices}

    \part[2] An algorithm consists of two phases executed sequentially. Phase 1 takes \(\Theta(n)\) time, and Phase 2 takes \(O(n^2)\) time. Which statements must be true about the total running time \(T(n)\)?
    \begin{choices}
        \choice \(T(n)=O(n^2)\).
        \choice \(T(n)=\Omega(n)\).
        \choice \(T(n)=\Theta(n^2)\).
        \choice \(T(n)=\Theta(n)\).
    \end{choices}

    \part[2] Which of the following postfix (Reverse Polish) expressions are \textbf{illegal}? All operators are binary, and division uses real-number arithmetic. A valid expression must never require more operands than are available and must leave exactly one value on the stack.
    \begin{choices}
        \choice \ttt{1 2 * + 3 5 +}
        \choice \ttt{4 5 6 / * 1 /}
        \choice \ttt{1 + 2 - 3 + 4}
        \choice \ttt{7 8 9 1 + - *}
    \end{choices}
\end{parts}
